\chapter{You've joined the J.A.P.I.E.-cie. Where to start?} \label{chap:getting-started/where-to-start}

Again, welcome to the committee! This chapter will guide you through this installation guide so that you can get your local development environment set up as quickly as possible.

This document is intentionally kept short and focused. If you are looking for broader background knowledge on the technologies used within the J.A.P.I.E.-cie, such as HTML, CSS, JavaScript, PHP, CakePHP, or Git workflows, please refer to the full J.A.P.I.E.-cie Reference.

\cref{chap:getting-started/terminology} contains a general list of terminology used within the J.A.P.I.E.-cie, mostly descriptions of the various software packages and technologies used by the committee. Use this chapter as a reference whenever you encounter a term that you are not familiar with.

\section*{Practicalities}

The committee uses a support ticket system to manage any problems or requests that the board---or sometimes, other association members---have for the committee. Details on this system are available in \cref{chap:getting-started/osticket}.

Many services and accounts are protected by a password. The system we use to manage these passwords is explained in \cref{chap:getting-started/passwords}.

\section*{Creating a local development environment}

The main goal of this document is to help you set up a local development environment so you can develop and test changes to the J.A.P.I.E.-cie websites on your own machine. Follow the instructions in \cref{chap:installation-guides/git,chap:installation-guides/local-web-server,chap:installation-guides/local-websites} in order to get everything set up.

\clearpage
